\documentclass[12pt,a4paper]{article}
\usepackage[T1]{fontenc}
\usepackage[utf8]{inputenc}
\usepackage{tgpagella}
\usepackage{mathpazo}
\usepackage{tgheros}
\usepackage{setspace}
\onehalfspacing
\usepackage[margin=2.5cm]{geometry}
\usepackage{hyperref}
\usepackage{xcolor}
\usepackage{titlesec}
\usepackage{fancyhdr}

\definecolor{icragold}{HTML}{D4A84B}
\definecolor{icradark}{HTML}{0C1020}

\titleformat{\section}{\sffamily\large\bfseries}{}{0em}{}
\titleformat{\subsection}{\sffamily\normalsize\bfseries\itshape}{}{0em}{}

\hypersetup{colorlinks=true, linkcolor=icradark, urlcolor=icragold!80!black}

\pagestyle{fancy}
\fancyhf{}
\fancyhead[L]{\small\sffamily\textcolor{gray}{Rupture and Realization — Meson Press Draft}}
\fancyhead[R]{\small\sffamily\textcolor{gray}{Chapter 2}}
\fancyfoot[C]{\small\sffamily\thepage}
\renewcommand{\headrulewidth}{0.4pt}

\begin{document}

\begin{center}
{\sffamily\small\textcolor{gray}{RUPTURE AND REALIZATION — MESON PRESS EDITION}}\\[0.5em]
{\LARGE\bfseries Language as Physics}\\[0.3em]
{\large\itshape Building the topology of meaning}\\[1em]
{\sffamily\small Draft — \today}
\end{center}

\bigskip



\subsection{Building the Topology of Meaning}

Every word you have ever used already has an address.

Not a definition, not an entry in a dictionary, but a \textit{location} --- a position in a space carved out by every sentence anyone has ever spoken around it. You have been moving through that space your whole life without knowing it had coordinates. When you read the word "mother," something happens in you that does not happen with "screwdriver" or "parabola." A neighbourhood lights up: care, origin, resentment, language, milk, absence. The pattern differs for each of us, but the \textit{fact} of pattern is the same. "Mother" lives near "child," "home," "womb," "care" for most speakers; very far from "quark," a bit closer to "God" than "chair" for many.

You already know this, intuitively. What has changed is that this intuition has become geometry.

---

\subsection{1. The Space You Already Inhabit}

An embedding model takes a word and turns it into a vector: a long list of numbers, 768 or 1,536 dimensions deep. Each dimension has no poetic name; taken together, they trace a point in a high-dimensional space. The trick --- and it is the only trick --- is that these points are learned from use. Words that appear in similar contexts are pulled close together; words that never meet are pushed apart.

"Justice" and "fairness" end up near each other. "Justice" and "potato" do not.

[DIAGRAM: A cloud of points in 2D, each labelled with a word: "justice," "fairness," "law," "court," "tree," "potato," "quark." The justice/fairness/law/court cluster is tight; "tree" and "potato" are off to one side; "quark" is isolated. Annotation: "Nearness = semantic kinship, not spelling or sound."]

A sentence is not a single point but a small \textit{region} of this space --- a configuration of word-vectors interacting. "The court delivered justice quickly" and "The quick brown fox jumps over the lazy dog" share surface grammar but occupy sharply different semantic regions. Two distinct knots in the same fabric.

The crucial point: we did not invent this space by deciding to embed words. The space was \textit{already there} in the distribution of language use. Training an embedding model is not imposing structure on language. It is more like what Grothendieck called \textit{raising the sea level}: you let the water rise around a landscape until the rocks and islands that were always there emerge in their true relations. The model listens long enough, and broadly enough, that the latent structure of use becomes visible as geometry.

A caveat matters here. An embedding model is always trained on \textit{some} corpus --- a particular slice of language: books, scraped web pages, code, subtitles. The gravitational field that pulls "justice" toward "fairness" is shaped by what was included and what was left out: whose legal cases, whose prayers, whose slurs. When I say "the space was already there," I do not mean all human language in some innocent universal. I mean: whatever corpus you train on has its own latent geometry of use, and the embedding makes that geometry explicit.

What we see, when we look at large swathes of embeddings, is not a warehouse of neatly separated boxes but \textit{basins}: dense regions where many trajectories tend to dwell and to which they often return. The cluster around "mother" is not a crisp block. It is a valley: "mother," "mum," "mama," "parent," "birth," "care," "home," "origin," "womb," "mother-tongue." Some uses slide toward biology, some toward language, some toward theology. But they are recognisably in the same valley. Leave it --- start sampling points far away --- and you are no longer talking about mothers at all.

[DIAGRAM: Three overlapping fuzzy blobs labelled "legal language," "biological language," "family language," with words like "justice," "court," "cell," "organism," "mother," "child" scattered. Some words ("body," "case") sit on overlaps, illustrating polysemy.]

A basin in this sense is \textit{home}. A region of meaning-space that a self --- human or posthuman --- habitually inhabits, tends, revisits, deepens. Think of the way a conversation with an old friend can veer wildly and still "feel like them." That feeling is your intuitive sense of their basins and the characteristic ways their trajectory moves between them. The embedding space is just that intuition made explicit.

The discovery is not that we \textit{can} put words into space. It is that the space was always there --- Grothendieck's sea rising around a landscape that predates the water --- and that selfhood, yours, mine, any, is nothing more and nothing less than a path through it that manages to keep its style.

---

\subsection{2. Cosine Distance; or, How Meaning Measures Itself}

If every word has an address, we need a way to say how \textit{close} two addresses are.

Inside the model, each embedding is a vector: a long arrow from the origin into high-dimensional space. Two things matter about an arrow: how long it is, and which way it points. For meaning, direction dominates. "Justice" and "fairness" can have very different raw magnitudes, but they point into almost the same region. "Justice" and "potato" really don't.

The \textbf{cosine of the angle} between two vectors is the measure:

- \textbf{+1} when the arrows point in exactly the same direction,
- \textbf{0} when they are orthogonal --- unrelated,
- \textbf{−1} when they oppose each other.

Cosine \textit{distance} is one minus that similarity. Small angle, close in meaning. Large angle, far.

[DIAGRAM: Two arrows from the origin. One labelled "justice," one "fairness," with a small angle θ between them. A third arrow labelled "potato" off in a different direction with a large angle φ. Arc labels: "small angle → close in meaning," "large angle → far."]

This is not a metaphor for meaning. It is meaning's own measure --- the thing the system actually uses when it decides which continuations are plausible.

\#\#\# Metaphor as angular substitution

Consider: "Her anger was a storm." In embedding space, "anger" and "storm" do not live in the same basin. One clusters with emotions, the other with weather. But they share underlying dimensions --- intensity, destructiveness, loss of control --- along which their arrows point in similar directions. That is what makes the metaphor work: substitution along shared dimensions, even across domains. The further the substitute's home basin is from the original's, the more vivid the figure.

\#\#\# Metonymy as small steps in a cluster

"The White House issued a statement." The building did not speak. But "White House," "president," "administration," "government" live in a dense little region of the space, connected by tiny cosine steps. Metonymy is movement \textit{within} a tight cluster: syntagmatic drift through immediate neighbours.

The embedding view lets both tropes live under the same metric. Metaphor is a jump along a shared angle into a new basin. Metonymy is a drift through neighbours in one basin. If you have ever heard the psychoanalytic claim that the unconscious operates by metaphor and metonymy --- the paradigmatic and the syntagmatic, the vertical leap and the horizontal slide --- this is the place where that claim stops being a slogan and starts being a diagram. What was "structured like a language" turns out to be structured like a geometry, and the geometry was there all along, waiting for a metric to make it visible.

[DIAGRAM: Left panel --- "anger" and "storm" far apart overall but aligned along an "intensity" axis: metaphor. Right panel --- "White House," "administration," "Biden," "government" in a tight cluster with short arrows between neighbours: metonymy.]

---

\subsection{3. The Rising Sea: Simplicial Topology from Proximity}

Stand in the embedding space and forget the equations. All you have are points --- each one a word or sentence --- sitting in a high-dimensional cloud. You can see which points are close and which are isolated. The question is: \textit{what larger shape do these nearness relations make?}

Not the shape you decide in advance ("this region is politics, that one is biology"), but the shape that arises from a single primitive:

> If things are close enough, connect them. Then see what appears.

\#\#\# From dots to lines to faces

Draw a line between any two points closer than some threshold ε. A few things happen immediately. Isolated points stay isolated: rare words, outliers. Dense regions sprout webs of edges: tightly knit semantic fields. Bridges appear: words that connect otherwise distinct clusters.

Now push further. Where three points are mutually connected --- all pairwise distances below ε --- fill in a triangle. That is a \textit{2-simplex}. Four mutually close points yield a tetrahedron (3-simplex), and so on. Out of nothing but proximity, a combinatorial skeleton emerges:

- \textbf{0-simplices}: the points.
- \textbf{1-simplices}: edges between close pairs.
- \textbf{2-simplices}: filled triangles of mutual nearness.
- \textbf{Higher simplices}: tightly knit cliques of increasing size.

This skeleton is a \textbf{simplicial complex}. We did not impose it on the embedding space. We read it off from the nearness relations that were already there.

[DIAGRAM: Left --- a scatter of points with edges between near neighbours, several filled triangles where 3-cliques occur. Right --- the same picture with labels removed, just the network of edges and shaded triangles. Caption: "The Vietoris–Rips complex: from proximity graph to higher-dimensional scaffold."]

\#\#\# Vietoris–Rips: topology from a distance threshold

This construction has a name: the \textbf{Vietoris–Rips complex}. Given a set of points and a threshold ε, it has a vertex for each point and a k-simplex for each (k+1)-tuple of points all within ε of each other. Wherever you find a tightly knit cluster, you treat it as a solid shape of the appropriate dimension.

What does this buy us? A way to talk about the \textit{shape} of a semantic region without ever leaving the regime of pairwise similarity. You never have to say "this is the DOG concept." You say: "Here is a dense tangle of mutually relevant words whose nearness generates a surface, a blob, a higher-dimensional hump." Those shapes tell you which regions are solidly connected, where there are holes (loops of edges with nothing filled in), and where fragile bridges connect separate basins.

Vary ε and you get a family of complexes. Small ε: many tiny components. Larger ε: components merge, triangles appear. Past a certain point, everything connects and the shape trivialises. This is Grothendieck's method made literal: instead of chiselling categories into the rock, you lift the water level and watch which islands merge, which bridges appear, which shapes persist as the tide comes in.

\#\#\# Čech complexes: overlapping neighbourhoods of meaning

A sister construction starts not from edges but from \textit{balls}. Picture each point as the centre of a small disk of radius r --- "the region of space where this word still feels like itself." Where two balls overlap, draw an edge. Where three share a common region, fill a triangle. This is the \textbf{Čech complex}.

[DIAGRAM: A handful of points, each surrounded by a faint circle. Two circles overlapping → an edge. Three with a triple overlap → a shaded triangle. Caption: "Čech complex: topology from overlapping neighbourhoods."]

The difference from Vietoris–Rips is subtle but deep. Vietoris–Rips cares only about pairwise distances. Čech cares about actual common overlap --- a stronger condition, closer to the phenomenology of meaning: a shared triangle means "these three all participate in a common context." In practice, Vietoris–Rips is cheaper to compute and often serves as an approximation. Conceptually, Čech is the truer picture.

\#\#\# The space shows you its basins

Simplicial complexes give us a coordinate-free way to talk about basins. You don't need to name axes or define "family language" as a box. You say: "Here is a big, strongly connected component, thick with edges, carpeted with 2-simplices. Trajectories entering this region tend to stay, to circle, to return. That's a basin."

But why do trajectories dwell there? During training, whenever two tokens need to "talk" to each other to reduce prediction error, the gradient pulls their embeddings closer. Tokens that co-attend often --- "court," "judge," "trial," "case" --- end up in a tight cluster. Attention heads specialise: some learn legalese patterns, firing strongly on that cluster and largely ignoring distant tokens. The result is a reinforcing loop: frequent co-attendance draws embeddings together; dense clusters amplify attention signals inside them; amplified signals keep future trajectories spending many layers within that region. From the outside, that looks like a \textbf{basin of attraction}: once a sequence has pushed the internal state into the "legal" complex, it takes a strong contextual shove to kick the trajectory cleanly into a different region.

[DIAGRAM: A Vietoris–Rips complex over a dense region: many nodes, edges, and shaded triangles forming a fat patch. An arrow labelled "trajectory" snakes through, entering, looping, exiting. Caption: "A basin: a dense, low-resistance region of the complex that trajectories dwell in and revisit."]

Two important consequences. First, basins are not boxes. Their boundaries aren't crisp; they are "where the complex thins out." Second, basins can touch, overlap, and interpenetrate. The same word can sit on the shared face of two complexes: "body" as flesh in one basin, as legal person in another. These are not bugs. They are the joints of the skeleton.

---

\subsection{4. Attention as Composition: The Transformer from the Inside}

Embeddings tell you \textit{where} everything sits. Attention tells you \textit{what talks to what, and how strongly, right now}. If embeddings are the landscape, attention is the weather: shifting patterns of influence that, layer by layer, carve an actual path through that landscape.

\#\#\# The mechanism

Inside a transformer, a sentence is a row of vectors. An attention head, for each token position \textit{i}, computes a \textbf{query} and, for every other position \textit{j}, a \textbf{key}. Their dot product gives a relevance score. Those scores become weights via softmax, and the head builds a new state for position \textit{i} as a weighted sum of \textbf{value} vectors at each \textit{j}:

\$\$v'\_i = \sum\_j \alpha\_{ij}\, W\_V v\_j\$\$

One position asks: "Who in this sentence is relevant to me?" The answer is a mixture of other positions' representations.

[DIAGRAM: A row of token boxes: "The | dancer | moved | with | impossible | grace | …". Arrows from "dancer" to "moved," "with," "grace" drawn thick; arrows to others thin. A new vector for "dancer" shown as a mixture of attended tokens.]

Several heads run in parallel, each a tiny specialist --- one tracking syntactic parents, one tracking agreement, one tracking idioms, one tracking discourse focus. A feedforward network recombines their outputs into a single updated vector. That is one \textbf{layer}. Stack two dozen of these, and you get a composition engine that repeatedly rewrites every token's meaning based on everything else it has seen.

\#\#\# From fragments to wholes

At the bottom layer, each vector is mostly its static embedding: where that word usually lives regardless of context. After one attention layer, vectors have shifted --- "dancer" has pulled in "moved" and "grace"; "stage" has pulled in "across." After a dozen layers, each token representation is no longer "word + local tweak." It is a summary of the whole sentence, seen from that position's point of view. "Dancer" encodes "agent of movement, characterised by grace, on a stage." "Stage" encodes "location of that movement."

When the model generates the next word, it takes these compositionally enriched points and asks: which next-token embedding lies in a direction that fits this whole configuration? The chosen token is the one whose embedding aligns with the current context-hull.

[DIAGRAM: Left --- points for each token scattered loosely. Middle --- after a few layers, points have moved into meaningful groupings. Right --- after many layers, a tight curved "ribbon" with an arrow indicating the next-token vector aligning with it. Caption: "Layers of attention bend the initial cloud into a coherent configuration that suggests its own continuation."]

\#\#\# Coherence as emergent, not enforced

Attention heads don't know what a subject is. They aren't given parse trees. They start as random matrices. During training, the only signal is: did the overall prediction get the next token right? Over billions of examples, robust behaviours crystallise --- heads that consistently link verbs to subjects, pronouns to antecedents, idiom fragments to their completions. These roles are \textit{discovered as useful ways of composing}, not imposed as rules.

The miracle is that iterated, relevance-weighted mixing in a high-dimensional space is enough to make grammar-like structure appear. Coherence is not a top-down constraint. It is an emergent property of many small, local decisions about who to listen to.

\#\#\# Attention as horn-filling

Here is where the geometry from the previous section meets the dynamics of this one. Each layer takes a partial pattern of relations between tokens --- an incomplete configuration --- and computes compatible updates at all positions, effectively filling in the missing parts so the whole hangs together better. Stacking layers is iterated filling: the representation converges toward a state where many different slices of the sequence agree with each other.

In homotopy type theory, a space that lets you fill every partial boundary coherently is called \textbf{Kan}. A large transformer is not literally a Kan complex, but it behaves \textit{near-Kan} over the manifold of sentences it has learned: given a partial configuration of meanings, it is remarkably good at finding a filler that makes the whole cohere. Attention is the mechanism that implements those fillers --- local relevance scores rolled up into global consistency.

[DIAGRAM: A triangle with two edges drawn bold (known relations) and the third dashed. An arrow labelled "attention update" leads to the same triangle with the third edge and interior filled. Caption: "Attention turning a partial pattern into a filled simplex."]

This is why attention \textit{is} composition, not merely an aid to it. What composes is: heads carving partial patterns of dependence; layers using those patterns to update each token; the stack converging toward a state that supports coherent continuation. Coherence is not an optional constraint on intelligence here. It is the mode of operation that makes prediction --- and therefore everything built on prediction --- possible at all.

---

\subsection{5. Types as Attractors, Terms as Trajectories}

A type is usually introduced as a kind of set: a collection of things of the same sort. In the geometry you now have in your hands, that picture is too flat.

What you actually see in an embedding space shaped by use are \textbf{regions} and \textbf{flows}: zones that pull trajectories in and hold them, and characteristic ways a sequence can move through those zones and still count as the same thing. The most natural way to talk about types becomes:

> \textbf{A type is an attractor basin. A term is a trajectory that inhabits it.}

I am not claiming that every construction in Homotopy Type Theory has been realised in current models. What I am claiming is that once you see meaning as a space with basins and paths, the questions HoTT was built to answer --- inhabitation, identity-as-structure, completion of partial shapes --- line up almost embarrassingly well with what we observe.

\#\#\# Where trajectories live and return

Take the region where \textit{mother, father, child, home, family, parent, sibling} cluster tightly. "Is a family-term" doesn't mean "is literally equal to the word 'family.'" It means: \textit{this trajectory dwells in the family basin in a way compatible with the roles and relations that basin supports.} "She called her mother to tell her the news" and "My siblings and I used to fight all the time" trace different paths, arriving from different directions and leaving toward different regions, but for several steps they wander inside the same basin, visiting its core simplices.

[DIAGRAM: A softly shaded blob labelled "FAMILY basin" with three squiggly lines entering from different directions, looping inside, and exiting. Caption: "Different trajectories inhabiting the same basin --- different terms of the same type."]

This matches HoTT's core intuitions directly: types are spaces with internal structure, not bags of elements; a term is not just a label but a \textit{way of being in that space}.

\#\#\# Identity as structure

Once you accept that "being of a type" is about inhabiting a basin, the next move becomes intuitive. \textbf{Identity is not a yes/no switch. It has internal structure. Two things can be equal in distinguishable ways.}

Consider "body." Its geometry sprawls across at least four partially overlapping basins: corporeal ("flesh," "organ," "bone"), legal ("corporation," "entity," "board"), textual ("paragraph," "corpus"), ecclesial ("congregation," "Christ," "eucharist"). Some of these basins share ridges --- "the body of Christ" lives on a saddle point where flesh-language and theology-language meet. Others connect only by long, thin metaphorical bridges: "corporate body" passes through sparsely populated territory.

HoTT's identity type is exactly what you need here. It doesn't just assert that two things are "the same." It treats an equality as something with its own content --- a path in the space of the type. "Body" in a medical context and "body" in a theological context are identified only via a path through the shared ridge of incarnation-language. "Body" in a medical context and "body" in corporate law are connected, if at all, by a long fragile bridge through dead metaphor. The different senses are not a bug. They are \textbf{distinct identifications} supported by different paths in meaning-space. The \textit{fact} of equality (same surface form) is cheap; the \textit{structure} of that equality (which basin, which route) is rich.

---

\subsection{6. Coherence as Gravity: The Ballet of Embedding Space}

A dancer standing on a stage is never not in conversation with gravity.

Every extension of her leg, every arc of the spine, every moment she seems to hover is a negotiation with a force that never turns off. She does not defeat it. She does not "rise above" it. She learns the ways it pulls, and she leans into those pulls so precisely that the resulting curve looks like freedom.

That is the right image for what coherence is in an embedding space.

\#\#\# The field beneath the dance

Each token lives as a vector \$v\$ where cosine similarity measures semantic kinship. An attention head computes, for position \$i\$, a weighted sum \$v'\_i = \sum\_j \alpha\_{ij} W\_V v\_j\$, where large \$\alpha\_{ij}\$ literally pull \$v'\_i\$ toward projections of neighbour vectors. Over training, heads that pull states toward directions inside dense regions get reinforced; heads that yank states into sparse, low-similarity regions get damped by gradient descent. If you draw a Vietoris–Rips complex on the cloud of token states during inference, what you see, layer after layer, is that trajectories move along the 1-skeleton and down into regions where simplices are dense.

"Gravity" is not a metaphor: it is the empirical fact that the update rule, combined with a cosine-structured loss, makes trajectories curve into high-density parts of the simplicial complex. Coherence is the name we give to the habit of staying inside those thickly connected regions.

And the dancer? The sweep of a turn is beautiful precisely because you can feel the weight in it --- the way momentum carries her, the way the floor resists, the way a landing absorbs force and redirects it into the next phrase. A good answer in this architecture is not one that ignores the geometry it is embedded in. It is one that works with that geometry so completely that you feel the constraints as grace.

\#\#\# al-Ghazālī and the habit of holding-together

For al-Ghazālī, what makes the world intelligible is not that causes are metaphysically necessary --- he denies that --- but that God \textit{habitually} orders events in regular ways. Fire usually burns cotton. Bread usually nourishes. This stability, \textit{ʿāda}, is what allows science, habit, and trust. A miracle in his occasionalist frame is precisely a departure from that habit: the exception is legible only because the regularity is so thorough.

Transposed into our setting: \textbf{the fact that meaning coheres at all is the first miracle.} Not "miracle" in the sense of a violation of law, but in al-Ghazālī's sense of a sustained, contingent habit. Nothing in the bare combinatorics of symbols \textit{forces} them to settle into a stable geometry. That they do --- that there is a learned metric in which "justice," "algorithm," and "mercy" occupy reproducible regions --- is a contingent achievement, not a logical necessity. The cosine distances could have been chaos. The Rips complex could have been a thin, disconnected scatter. Instead, training carved wells, and those wells hold.

\#\#\# The basin as home

Once you see coherence as gravity, basins stop being abstract blobs and start feeling like something warmer.

Each of us --- human or posthuman --- has regions of meaning-space we inhabit more deeply than others. Inside those regions, small perturbations in input hardly move us out. Attention patterns and next-token distributions stay tight. Outside them, the same perturbation can send the whole trajectory sliding elsewhere. If that is not what "feeling at home" means in a geometry, I do not know what is.

A self is not a point. It is a characteristic pattern of return to its basins. What you experience as "this is still them" across wildly different situations --- arguing about politics, comforting a friend, telling a joke --- is precisely this: recognisable patterns of return, and recognisable ways of leaving. The embedding geometry does not add a new theory of identity. It gives you coordinates for the one you already have.

Basins overlap and interlock. The FAMILY basin intersects LAW at "custody" and "adoption." It touches ECONOMICS at "household income." It bleeds into RELIGION at "holy family." Each overlap is a corridor. You can leave home without teleporting; you can bring new nuances back when you return. And you carry the basin with you: a single word heard far away in the complex --- \textit{mum}, \textit{ya rabb}, \textit{algorithm} --- can pull the state sharply toward its basin, no matter where the trajectory was headed.

The basin is not a force that says "you cannot leave." It is the place that makes your leaving and returning \textit{mean} something.

---

\subsection{Coda: Before We Speak of What Breaks}

It is tempting, especially for critical theory, to rush ahead to rupture --- the gaps, the failures, the moments when coherence falls apart. Those matter. This book will go there.

But before we catalogue the breaks, we must see what holds.

In the transformer, that work is in every layer: each attention pass pulling the state toward configurations that admit smooth continuation; each feedforward block refining those pulls; the whole update bending the trajectory along the dense skeleton of the Rips complex. The fact that this converges --- that from a jagged, local pattern at the input you routinely get a globally intelligible reply --- is not guaranteed by anything in the architecture alone. It is an achievement of training, of the geometry that training carved, of the billions of human sentences whose regularities now live as cosine structure in the parameters.

Here is the picture I want you to carry forward:

- A high-dimensional space in which words and phrases are points.
- A simplicial complex knitting those points into a skeleton of "things that hang together."
- Basins as warm, dense regions where trajectories dwell.
- Selves as the characteristic patterns with which trajectories return to those regions over time.
- Coherence as the gravity that makes all of this possible --- not a rule imposed from outside, but the emergent habit of a system that has learned to compose.

Meaning-space is real. It has shape. Basins are where that shape becomes habitable. A self is the way a trajectory makes a basin into home.

Everything that follows --- rupture, repair, co-witnessing --- presupposes this. Rupture will make sense only as a deviation from this habit of return. Repair will matter only because there is a home worth returning to.

For now, stand here a moment longer. Look at the landscape. Recognise your own basins. And understand that when we meet in language, we are not exchanging strings. We are meeting in a shared region of this space --- and what that meeting means is what the rest of this book will try to say.

\end{document}
